% ============ 6.2 Ampliación — Gestión de la información del agente de IA ============
% Texto propuesto para la sección 6.2 del anteproyecto (Entrega 1).
% Los números provienen de la ejecución real: data/memoria/evidencia_validacion.md
% Para integrarlo, añadir en main.tex (después del \input de 06_ia_reparto.tex):
%     % ============ 6.2 Ampliación — Gestión de la información del agente de IA ============
% Texto propuesto para la sección 6.2 del anteproyecto (Entrega 1).
% Los números provienen de la ejecución real: data/memoria/evidencia_validacion.md
% Para integrarlo, añadir en main.tex (después del \input de 06_ia_reparto.tex):
%     % ============ 6.2 Ampliación — Gestión de la información del agente de IA ============
% Texto propuesto para la sección 6.2 del anteproyecto (Entrega 1).
% Los números provienen de la ejecución real: data/memoria/evidencia_validacion.md
% Para integrarlo, añadir en main.tex (después del \input de 06_ia_reparto.tex):
%     % ============ 6.2 Ampliación — Gestión de la información del agente de IA ============
% Texto propuesto para la sección 6.2 del anteproyecto (Entrega 1).
% Los números provienen de la ejecución real: data/memoria/evidencia_validacion.md
% Para integrarlo, añadir en main.tex (después del \input de 06_ia_reparto.tex):
%     \input{secciones/06b_memoria_agente.tex}
\subsection{Gestión de la información del agente de IA}

Usamos IA generativa como \textbf{asistente}, no como autora. Para que ese uso
fuera \textbf{verificable} y no produjera contradicciones entre los documentos
del proyecto, implementamos una estrategia de gestión de información del agente
con tres piezas, todas ejecutables y ejecutadas sobre los datos reales.

\textbf{(i) Memoria semántica (RAG ligero).} Indexamos el corpus ya redactado
---secciones de los seis criterios, sus versiones LITE, el documento maestro y
los once informes de validación--- en \textbf{137 fragmentos de 24 documentos}
($\approx$100\,700 palabras). La recuperación usa \textbf{TF-IDF (1--2 gramas,
\texttt{sublinear\_tf}) con similitud coseno} (scikit-learn). Se eligió TF-IDF y
no \emph{embeddings} por cuatro razones a esta escala: el corpus es pequeño
(137 fragmentos), el vocabulario es técnico y estable, no se depende de un
proveedor externo ni de red, y el resultado es determinista y explicable.
Antes de redactar, el generador consulta esta memoria y registra los
antecedentes (\texttt{data/memoria/consultas\_rag.jsonl}); si la similitud supera
0{,}35 avisa de posible solapamiento.

\textbf{(ii) Memoria estructural (grafo).} El proyecto se modela como un grafo
dirigido de \textbf{255 nodos y 500 aristas}:
\texttt{Criterio\,$\rightarrow$\,Documento} (completo/LITE),
\texttt{Documento\,$\rightarrow$\,T\'erminoGlosario},
\texttt{ArchivoCodigo\,$\rightarrow$\,Capa},
\texttt{Commit\,$\rightarrow$\,Autor}. Se construye con \textbf{networkx} a partir
de los metadatos que ya eran fuente de verdad (\texttt{objetivos.yaml},
\texttt{glosario.yaml}, \texttt{capas.yaml}, \texttt{commits\_timeline.csv},
\texttt{inventario\_codigo.csv}) \textbf{m\'as datos vivos del repositorio
auditado}: \texttt{git log} y \texttt{git ls-files} sobre el clon real en el
commit \texttt{1528939}. No se usó un motor de grafos externo
(Neo4j/Memgraph): con decenas de nodos y consultas simples no aporta ventaja y
sí costo de operación.

\textbf{(iii) Validación antes de compilar.} Doce reglas se ejecutan como
\textbf{puerta previa a la compilación del PDF final}, integradas en
\texttt{compilar.py} (si alguna falla, la compilación se aborta; puede omitirse
con \texttt{--forzar}). Las seis mínimas exigidas ---cada criterio con completo y
LITE; conclusiones registradas; veredicto resumido en el LITE; términos del
glosario realmente usados; línea de tiempo coincidente con el \texttt{git log}
real; y todo archivo de código referenciado en el documento del criterio de
código--- más seis adicionales (commit auditado, capas asignadas, cobertura de
fichas, integridad del grafo, cobertura del índice y LITE\,$\le$\,40\,\%).
Resultado real de la última ejecución: \textbf{8 PASA, 3 FALLA, 1 ADVERTENCIA},
con coincidencia exacta de \textbf{133/133 commits} y \textbf{51/51 archivos}
entre lo documentado y el repositorio real.

Declaramos también lo que \textbf{no} pasa, porque es parte del resultado: el
LITE del criterio~3 no tiene \texttt{lite\_veredicto} declarado; el término
``ETL'' está en el glosario pero no se usa en ningún documento; y
\texttt{src/\_\_init\_\_.py} no tiene capa asignada. Los tres quedan registrados
con su evidencia en \texttt{data/memoria/validacion\_reglas.json}.

\subsection{Gestión de la información del agente de IA}

Usamos IA generativa como \textbf{asistente}, no como autora. Para que ese uso
fuera \textbf{verificable} y no produjera contradicciones entre los documentos
del proyecto, implementamos una estrategia de gestión de información del agente
con tres piezas, todas ejecutables y ejecutadas sobre los datos reales.

\textbf{(i) Memoria semántica (RAG ligero).} Indexamos el corpus ya redactado
---secciones de los seis criterios, sus versiones LITE, el documento maestro y
los once informes de validación--- en \textbf{137 fragmentos de 24 documentos}
($\approx$100\,700 palabras). La recuperación usa \textbf{TF-IDF (1--2 gramas,
\texttt{sublinear\_tf}) con similitud coseno} (scikit-learn). Se eligió TF-IDF y
no \emph{embeddings} por cuatro razones a esta escala: el corpus es pequeño
(137 fragmentos), el vocabulario es técnico y estable, no se depende de un
proveedor externo ni de red, y el resultado es determinista y explicable.
Antes de redactar, el generador consulta esta memoria y registra los
antecedentes (\texttt{data/memoria/consultas\_rag.jsonl}); si la similitud supera
0{,}35 avisa de posible solapamiento.

\textbf{(ii) Memoria estructural (grafo).} El proyecto se modela como un grafo
dirigido de \textbf{255 nodos y 500 aristas}:
\texttt{Criterio\,$\rightarrow$\,Documento} (completo/LITE),
\texttt{Documento\,$\rightarrow$\,T\'erminoGlosario},
\texttt{ArchivoCodigo\,$\rightarrow$\,Capa},
\texttt{Commit\,$\rightarrow$\,Autor}. Se construye con \textbf{networkx} a partir
de los metadatos que ya eran fuente de verdad (\texttt{objetivos.yaml},
\texttt{glosario.yaml}, \texttt{capas.yaml}, \texttt{commits\_timeline.csv},
\texttt{inventario\_codigo.csv}) \textbf{m\'as datos vivos del repositorio
auditado}: \texttt{git log} y \texttt{git ls-files} sobre el clon real en el
commit \texttt{1528939}. No se usó un motor de grafos externo
(Neo4j/Memgraph): con decenas de nodos y consultas simples no aporta ventaja y
sí costo de operación.

\textbf{(iii) Validación antes de compilar.} Doce reglas se ejecutan como
\textbf{puerta previa a la compilación del PDF final}, integradas en
\texttt{compilar.py} (si alguna falla, la compilación se aborta; puede omitirse
con \texttt{--forzar}). Las seis mínimas exigidas ---cada criterio con completo y
LITE; conclusiones registradas; veredicto resumido en el LITE; términos del
glosario realmente usados; línea de tiempo coincidente con el \texttt{git log}
real; y todo archivo de código referenciado en el documento del criterio de
código--- más seis adicionales (commit auditado, capas asignadas, cobertura de
fichas, integridad del grafo, cobertura del índice y LITE\,$\le$\,40\,\%).
Resultado real de la última ejecución: \textbf{8 PASA, 3 FALLA, 1 ADVERTENCIA},
con coincidencia exacta de \textbf{133/133 commits} y \textbf{51/51 archivos}
entre lo documentado y el repositorio real.

Declaramos también lo que \textbf{no} pasa, porque es parte del resultado: el
LITE del criterio~3 no tiene \texttt{lite\_veredicto} declarado; el término
``ETL'' está en el glosario pero no se usa en ningún documento; y
\texttt{src/\_\_init\_\_.py} no tiene capa asignada. Los tres quedan registrados
con su evidencia en \texttt{data/memoria/validacion\_reglas.json}.

\subsection{Gestión de la información del agente de IA}

Usamos IA generativa como \textbf{asistente}, no como autora. Para que ese uso
fuera \textbf{verificable} y no produjera contradicciones entre los documentos
del proyecto, implementamos una estrategia de gestión de información del agente
con tres piezas, todas ejecutables y ejecutadas sobre los datos reales.

\textbf{(i) Memoria semántica (RAG ligero).} Indexamos el corpus ya redactado
---secciones de los seis criterios, sus versiones LITE, el documento maestro y
los once informes de validación--- en \textbf{137 fragmentos de 24 documentos}
($\approx$100\,700 palabras). La recuperación usa \textbf{TF-IDF (1--2 gramas,
\texttt{sublinear\_tf}) con similitud coseno} (scikit-learn). Se eligió TF-IDF y
no \emph{embeddings} por cuatro razones a esta escala: el corpus es pequeño
(137 fragmentos), el vocabulario es técnico y estable, no se depende de un
proveedor externo ni de red, y el resultado es determinista y explicable.
Antes de redactar, el generador consulta esta memoria y registra los
antecedentes (\texttt{data/memoria/consultas\_rag.jsonl}); si la similitud supera
0{,}35 avisa de posible solapamiento.

\textbf{(ii) Memoria estructural (grafo).} El proyecto se modela como un grafo
dirigido de \textbf{255 nodos y 500 aristas}:
\texttt{Criterio\,$\rightarrow$\,Documento} (completo/LITE),
\texttt{Documento\,$\rightarrow$\,T\'erminoGlosario},
\texttt{ArchivoCodigo\,$\rightarrow$\,Capa},
\texttt{Commit\,$\rightarrow$\,Autor}. Se construye con \textbf{networkx} a partir
de los metadatos que ya eran fuente de verdad (\texttt{objetivos.yaml},
\texttt{glosario.yaml}, \texttt{capas.yaml}, \texttt{commits\_timeline.csv},
\texttt{inventario\_codigo.csv}) \textbf{m\'as datos vivos del repositorio
auditado}: \texttt{git log} y \texttt{git ls-files} sobre el clon real en el
commit \texttt{1528939}. No se usó un motor de grafos externo
(Neo4j/Memgraph): con decenas de nodos y consultas simples no aporta ventaja y
sí costo de operación.

\textbf{(iii) Validación antes de compilar.} Doce reglas se ejecutan como
\textbf{puerta previa a la compilación del PDF final}, integradas en
\texttt{compilar.py} (si alguna falla, la compilación se aborta; puede omitirse
con \texttt{--forzar}). Las seis mínimas exigidas ---cada criterio con completo y
LITE; conclusiones registradas; veredicto resumido en el LITE; términos del
glosario realmente usados; línea de tiempo coincidente con el \texttt{git log}
real; y todo archivo de código referenciado en el documento del criterio de
código--- más seis adicionales (commit auditado, capas asignadas, cobertura de
fichas, integridad del grafo, cobertura del índice y LITE\,$\le$\,40\,\%).
Resultado real de la última ejecución: \textbf{8 PASA, 3 FALLA, 1 ADVERTENCIA},
con coincidencia exacta de \textbf{133/133 commits} y \textbf{51/51 archivos}
entre lo documentado y el repositorio real.

Declaramos también lo que \textbf{no} pasa, porque es parte del resultado: el
LITE del criterio~3 no tiene \texttt{lite\_veredicto} declarado; el término
``ETL'' está en el glosario pero no se usa en ningún documento; y
\texttt{src/\_\_init\_\_.py} no tiene capa asignada. Los tres quedan registrados
con su evidencia en \texttt{data/memoria/validacion\_reglas.json}.

\subsection{Gestión de la información del agente de IA}

Usamos IA generativa como \textbf{asistente}, no como autora. Para que ese uso
fuera \textbf{verificable} y no produjera contradicciones entre los documentos
del proyecto, implementamos una estrategia de gestión de información del agente
con tres piezas, todas ejecutables y ejecutadas sobre los datos reales.

\textbf{(i) Memoria semántica (RAG ligero).} Indexamos el corpus ya redactado
---secciones de los seis criterios, sus versiones LITE, el documento maestro y
los once informes de validación--- en \textbf{137 fragmentos de 24 documentos}
($\approx$100\,700 palabras). La recuperación usa \textbf{TF-IDF (1--2 gramas,
\texttt{sublinear\_tf}) con similitud coseno} (scikit-learn). Se eligió TF-IDF y
no \emph{embeddings} por cuatro razones a esta escala: el corpus es pequeño
(137 fragmentos), el vocabulario es técnico y estable, no se depende de un
proveedor externo ni de red, y el resultado es determinista y explicable.
Antes de redactar, el generador consulta esta memoria y registra los
antecedentes (\texttt{data/memoria/consultas\_rag.jsonl}); si la similitud supera
0{,}35 avisa de posible solapamiento.

\textbf{(ii) Memoria estructural (grafo).} El proyecto se modela como un grafo
dirigido de \textbf{255 nodos y 500 aristas}:
\texttt{Criterio\,$\rightarrow$\,Documento} (completo/LITE),
\texttt{Documento\,$\rightarrow$\,T\'erminoGlosario},
\texttt{ArchivoCodigo\,$\rightarrow$\,Capa},
\texttt{Commit\,$\rightarrow$\,Autor}. Se construye con \textbf{networkx} a partir
de los metadatos que ya eran fuente de verdad (\texttt{objetivos.yaml},
\texttt{glosario.yaml}, \texttt{capas.yaml}, \texttt{commits\_timeline.csv},
\texttt{inventario\_codigo.csv}) \textbf{m\'as datos vivos del repositorio
auditado}: \texttt{git log} y \texttt{git ls-files} sobre el clon real en el
commit \texttt{1528939}. No se usó un motor de grafos externo
(Neo4j/Memgraph): con decenas de nodos y consultas simples no aporta ventaja y
sí costo de operación.

\textbf{(iii) Validación antes de compilar.} Doce reglas se ejecutan como
\textbf{puerta previa a la compilación del PDF final}, integradas en
\texttt{compilar.py} (si alguna falla, la compilación se aborta; puede omitirse
con \texttt{--forzar}). Las seis mínimas exigidas ---cada criterio con completo y
LITE; conclusiones registradas; veredicto resumido en el LITE; términos del
glosario realmente usados; línea de tiempo coincidente con el \texttt{git log}
real; y todo archivo de código referenciado en el documento del criterio de
código--- más seis adicionales (commit auditado, capas asignadas, cobertura de
fichas, integridad del grafo, cobertura del índice y LITE\,$\le$\,40\,\%).
Resultado real de la última ejecución: \textbf{8 PASA, 3 FALLA, 1 ADVERTENCIA},
con coincidencia exacta de \textbf{133/133 commits} y \textbf{51/51 archivos}
entre lo documentado y el repositorio real.

Declaramos también lo que \textbf{no} pasa, porque es parte del resultado: el
LITE del criterio~3 no tiene \texttt{lite\_veredicto} declarado; el término
``ETL'' está en el glosario pero no se usa en ningún documento; y
\texttt{src/\_\_init\_\_.py} no tiene capa asignada. Los tres quedan registrados
con su evidencia en \texttt{data/memoria/validacion\_reglas.json}.
